\documentclass[luatex,a4paper]{ltjsarticle}
\usepackage{luatexja-fontspec}
\setmainjfont{Noto Serif CJK JP}
\setsansjfont{Noto Sans CJK JP}
\usepackage{amsmath,amssymb,amsthm}
\usepackage{mathtools}
\usepackage{tikz-cd}
\usepackage{enumitem}
\usepackage{hyperref}
\hypersetup{
  colorlinks=true,
  linkcolor=blue,
  urlcolor=teal,
  citecolor=magenta
}

\title{構造塔の普遍性と最小層選択\\--- Structure Tower 発展課題 ST-CAT-002 解説ノート ---}
\author{Codex (OpenAI)\thanks{本資料（TeX）および対応する Lean ファイルは Codex (OpenAI) が生成しました。}}
\date{\today}

\theoremstyle{definition}
\newtheorem{definition}{定義}[section]
\newtheorem{example}[definition]{例}
\theoremstyle{plain}
\newtheorem{theorem}[definition]{定理}
\newtheorem{proposition}[definition]{命題}
\theoremstyle{remark}
\newtheorem{remark}[definition]{注意}

\begin{document}
\maketitle

\section{概要}
本ノートは Structure Tower 発展課題 ST-CAT-002 の改訂版 Lean ファイル
\texttt{CAT2\_revised.lean}（Codex (OpenAI) により生成）をもとに、
学部生を対象として構造塔の三つのバージョンを解説する。
特に Version~A における「最小層選択」による普遍性の完全な一意性を強調し、
Version~B, C との違いを理解できることを目標とする。

\section{構造塔の基本概念}
\subsection{層と被覆}
\begin{definition}
型 $T$ の\textbf{構造塔}とは、次のデータからなる。
\begin{enumerate}[label=(\arabic*)]
  \item 基礎集合（または型） $\mathrm{carrier}$。
  \item 層を添字付ける集合 $\mathrm{Index}$。
  \item 添字集合の前順序 $\leq$。
  \item 各添字に対応する部分集合 $\mathrm{layer}(i) \subseteq \mathrm{carrier}$。
\end{enumerate}
これらは「被覆性」
$\forall x,\, \exists i,\, x \in \mathrm{layer}(i)$ と
「単調性」
$i \leq j \Rightarrow \mathrm{layer}(i) \subseteq \mathrm{layer}(j)$
を満たす。
\end{definition}

学部生の直感としては、層は「段階的に情報が蓄積する集合列」を捉える。
ただし同じ要素が複数の層に同時に属してもよいため、
どの層が「一番小さい」かは別途議論が必要になる。

\subsection{射の考え方}
二つの構造塔 $T, T'$ の間の射 $\varphi : T \to T'$ は、
基礎写像 $\varphi_{\mathrm{map}} : T.\mathrm{carrier} \to T'.\mathrm{carrier}$ と
添字写像 $\varphi_{\mathrm{index}} : T.\mathrm{Index} \to T'.\mathrm{Index}$ の組である。
単調性と層の保存条件が課されるため、
「層構造を壊さない」写像として理解できる。

\section{Version A: 最小層を備える構造塔}
\subsection{最小層選択}
\begin{definition}
各要素 $x$ に対し「$x$ を含む最小の層」を返す関数
$\mathrm{minLayer} : T.\mathrm{carrier} \to T.\mathrm{Index}$ が存在していて、
レーマ
\[
  x \in \mathrm{layer}(\mathrm{minLayer}(x)),
  \qquad
  x \in \mathrm{layer}(i) \Rightarrow \mathrm{minLayer}(x) \leq i
\]
が成立するとき、$T$ を\textbf{最小層付き構造塔}と呼ぶ。
\end{definition}

Lean 実装では \texttt{StructureTowerWithMin} として定義される。
最小層は順序理論的な帰納的選択ではなく、
「階層が良い形で並んでいる」という仮定を明示する手段である。

\subsection{自由構造塔と普遍性}
前順序集合 $X$ に対して自由構造塔 $\mathrm{freeStructureTowerMin}(X)$ を定義する。
層は $\{\,x' \in X \mid x' \leq i\,\}$ で与えられ、最小層は恒等写像
$\mathrm{minLayer}(x) = x$ である。

\begin{theorem}
$f : X \to T.\mathrm{carrier}$ を任意の写像とし、
合成 $\mathrm{minLayer} \circ f$ が単調であると仮定する。
このとき、自由構造塔から $T$ への射で基礎写像が $f$ に一致するものが
一意的に存在する。
\end{theorem}
\begin{proof}[証明のアイデア]
最小層の情報を添字写像に用いると、
層の保存条件と単調性が自動的に満たされる。
任意の射 $\varphi$ が $f$ を基礎写像として持つとき、
最小層の性質から添字写像も一致することが分かり、一意性が導かれる。
\end{proof}

\paragraph{図式での理解}
\begin{center}
\begin{tikzcd}[row sep=large, column sep=huge]
X \arrow[d, "i_X"'] \arrow[r, "f"] & T.\mathrm{carrier} \arrow[d, "\mathrm{minLayer}"] \\
X \times X \arrow[r, "\mathrm{id} \times \mathrm{minLayer} \circ f"'] &
T.\mathrm{Index}
\end{tikzcd}
\end{center}
ここで $i_X$ は $x \mapsto (x,x)$ と書ける埋め込みであり、
右下の写像が単調であることで図式が可換になる。

\section{Version B: 一般的な構造塔}
\subsection{最小層を仮定しない}
最小層を備えない一般的な構造塔では、
射の添字写像が一意に決まらない場合がある。
Lean の定義 \texttt{StructureTower} では
$\mathrm{minLayer}$ フィールドが存在しない。

\subsection{自由構造塔と弱い一意性}
自由構造塔 $\mathrm{freeStructureTower}(X)$ は
層を一点集合 $\{i\}$ とし、添字の前順序を離散化することで構築される。
\begin{proposition}
任意の写像 $f : X \to T.\mathrm{carrier}$ に対して
射 $\varphi$ は存在するが、一意性は
基礎写像に限られる。すなわち
$\varphi_{\mathrm{map}} = \psi_{\mathrm{map}}$ は示せるが、
添字写像の一致は保証されない。
\end{proposition}

\subsection{積の普遍性（基礎写像レベル）}
\begin{theorem}
任意の二つの射 $f_1 : T \to T_1$, $f_2 : T \to T_2$ に対し、
直積構造塔 $T_1 \times T_2$ への射 $\langle f_1,f_2 \rangle$ が存在する。
さらに、射影との合成で元の射に戻る可換図式が成立する。
\end{theorem}
\begin{remark}
ただしここでも、射 $\langle f_1,f_2 \rangle$ の添字写像は必ずしも一意とは限らない。
\end{remark}

\section{Version C: 分離された構造塔}
\begin{definition}
構造塔 $T$ が\textbf{分離されている}とは、
異なる添字 $i \neq j$ に対し $\mathrm{layer}(i)$ と $\mathrm{layer}(j)$ が互いに素、
すなわち $\mathrm{layer}(i) \cap \mathrm{layer}(j) = \emptyset$ であることをいう。
\end{definition}
このとき各要素はただ一つの層に属するため、
Version~A の最小層を暗黙に持つ。
しかし前提が強すぎる場合も多く、実用性は限定的である。

\section{Version の比較}
\begin{center}
\setlength{\tabcolsep}{6pt}
\begin{tabular}{lccc}
\hline
特徴 & Version A & Version B & Version C \\
\hline
最小層 & 必須 & 不要 & 自動的に一意 \\
普遍性 & 射そのものが一意 & 基礎写像のみ一意 & 射そのものが一意 \\
一般性 & 中程度 & 高い & 低い \\
実装の複雑さ & 中程度 & 高い (添字の選択が必要) & 低い \\
\hline
\end{tabular}
\end{center}

\section{Lean 実装で学ぶポイント}
\begin{itemize}
  \item \textbf{構造体による抽象化}: Lean の \texttt{structure} を用いることで、
    数学的定義をそのままコード化できる。
  \item \textbf{普遍性の形式化}: 「存在し一意」という性質は
    Lean では \texttt{\textbackslash exists!} や
    \texttt{Equiv} などで記述される。
  \item \textbf{モノトン性の扱い}: 添字写像の単調性は
    層の単調性を引き継ぐ形で証明できる。
  \item \textbf{実装上の注意}: Version~B では
    選択関数の扱いが難しく、Lean でも明示的な選択が必要となる。
\end{itemize}

\section{学習のための演習問題}
\begin{enumerate}[label=\textbf{演習 \arabic*.}]
  \item Version~A の自由構造塔において、
    写像 $f : X \to T.\mathrm{carrier}$ が単調でない具体例を挙げ、
    普遍性が失敗する様子を説明せよ。
  \item Version~B の構造塔で、同じ基礎写像を持ちながら
    添字写像が異なる二つの射を具体的に構成せよ。
  \item Version~C の分離条件が壊れる例を作り、
    そのとき自然な最小層が定義できない理由を考察せよ。
\end{enumerate}

\section{まとめ}
構造塔の普遍性は「何を一意とみなすか」に依存する。
Version~A の最小層仮定は実装をやや制限するものの、
完全な一意性を保証し教育的価値が高い。
Version~B は一般的だが、実装と証明で選択関数に注意が必要である。
Version~C は特別な状況では簡潔に扱えるが、一般性を欠く。

今回の Lean 実装と本ノートを併読することで、
抽象概念が計算機上でどのように形式化されるか、
そしてその際に必要となる追加条件が何かを理解してほしい。

\end{document}
